\documentclass[a4paper,13pt,3p,oneside]{report}
\usepackage{tikz}
\usetikzlibrary{shapes.geometric, arrows.meta, positioning}
\usepackage{scrextend}
\changefontsizes{13pt}
\usepackage[utf8]{vietnam}
\usepackage[top=2cm,bottom=2cm,left=3.5cm,right=2.5cm]{geometry}
\usepackage{graphicx,fancybox,indentfirst,amsmath,amssymb,times,array,enumitem,titlesec,titletoc,chngcntr,multirow,fancyhdr,float,booktabs,xcolor,setspace,parskip,tabularx,longtable,hyperref}
\usepackage[ruled,vlined]{algorithm2e}
\usepackage[font=small,labelfont=bf]{caption}
\usepackage{listings}
\usepackage[all]{hypcap}
\usepackage{seqsplit}
\usepackage{makecell}
\usepackage{xurl}
\usepackage{hyperref}
\usepackage{tabularx}
\usepackage{ltablex}
\keepXColumns
\newcommand{\code}[1]{\begingroup\ttfamily\seqsplit{#1}\endgroup}
\hypersetup{pdfborder={0 0 0}}
\definecolor{backcolour}{rgb}{0.95,0.95,0.92}
\definecolor{codegreen}{rgb}{0,0.6,0}
\definecolor{codegray}{rgb}{0.5,0.5,0.5}
\definecolor{codepurple}{rgb}{0.58,0,0.82}
\lstdefinestyle{mystyle}{backgroundcolor=\color{backcolour},commentstyle=\color{codegreen},keywordstyle=\color{magenta},numberstyle=\tiny\color{codegray},stringstyle=\color{codepurple},basicstyle=\ttfamily\footnotesize,breaklines=true,captionpos=b,keepspaces=true,numbers=left,numbersep=5pt,showspaces=false,showstringspaces=false,showtabs=false,tabsize=2}
\lstset{style=mystyle}
\fancypagestyle{plain}{\fancyhf{}\fancyfoot[R]{\thepage}\renewcommand{\headrulewidth}{0pt}\renewcommand{\footrulewidth}{0pt}}
\setlength{\headheight}{18pt}
\titleformat{\chapter}[hang]{\centering\bfseries}{CHƯƠNG \thechapter.\ }{0pt}{}[]
\titlespacing*{\chapter}{0pt}{-20pt}{20pt}
\titleformat{\section}[hang]{\bfseries}{\thechapter.\arabic{section}\ \ \ \ }{0pt}{}[]
\titlespacing{\section}{0pt}{\parskip}{0.5\parskip}
\titleformat{\subsection}[hang]{\bfseries}{\thechapter.\arabic{section}.\arabic{subsection}\ \ \ \ }{0pt}{}[]
\titlespacing{\subsection}{30pt}{\parskip}{0.5\parskip}
\counterwithin{figure}{chapter}
\counterwithin{table}{chapter}
\setcounter{secnumdepth}{3}
\onehalfspacing
\setlength{\parskip}{6pt}
\setlength{\parindent}{15pt}
\def\TITLE{Xây Dựng Hệ Thống IoT Theo Dõi Sức Khỏe Và Trợ Lý Y Tế AI Có Xét Đến Kiến Trúc Và Bảo Mật}

\begin{document}
\begin{titlepage}
\thispagestyle{empty}
\begin{center}

{\textbf{\large{ĐẠI HỌC BÁCH KHOA HÀ NỘI}}}\\[0.3cm]

\vspace{3.0cm}
{\textbf{\huge{BÁO CÁO BÀI TẬP LỚN}}}\\[0.6cm]
\vspace{2.0cm}
{\textbf{\Large{Xây dựng hệ thống IoT giám sát chỉ số sinh hiệu thời gian thực dựa trên kiến trúc tự phòng thủ}}}\\[1.0cm]
\vspace{1.0cm}

{\textbf{\large{ĐÀM VIỆT ANH}}}\\
{\large{20252015M}}\\[0.15cm]
{\textbf{\large{VŨ THÙY LINH}}}\\
{\large{20252035M}}\\[0.15cm]
{\textbf{\large{NGUYỄN LÊ HÙNG}}}\\
{\large{20252036M}}\\[0.15cm]
{\textbf{\large{NGUYỄN THÙY DƯƠNG}}}\\
{\large{20252022M}}\\[0.4cm]

\vspace{0.6cm}
\begin{table}[H]
\centering
\resizebox{\textwidth}{!}{%
\begin{tabular}{ll}
\multicolumn{1}{c}{\textbf{Giảng viên hướng dẫn:}} & TS. Trịnh Văn Chiến \hspace{0.5cm} \\
                                                    & TS. Phạm Ngọc Hưng \hspace{0.5cm} \\
                                                    & PGS.TS. Trần Quang Đức \hspace{0.5cm} \\[0.25cm]
\multicolumn{1}{l}{\textbf{Trường:}} & Công nghệ thông tin và Truyền thông \hspace{0.5cm} \\[1cm]
\multicolumn{2}{c}{\textbf{HÀ NỘI, 07/2026}}
\end{tabular}%
}
\end{table}

\end{center}
\end{titlepage}

\newpage
\pagenumbering{roman}
\chapter*{TÓM TẮT NỘI DUNG}
Trong bối cảnh các thiết bị đeo và hệ thống chăm sóc sức khỏe từ xa ngày càng phổ biến, việc xây dựng một kiến trúc IoT y tế an toàn, có khả năng cập nhật dữ liệu thời gian thực và hỗ trợ người dùng diễn giải tình trạng sức khỏe là một yêu cầu thực tiễn. Đề tài này thực hiện xây dựng nguyên mẫu hệ thống IoT theo dõi sức khỏe sử dụng thiết bị ESP32-S3 kết nối cảm biến MAX30102 và DHT22 để đo nhịp tim, SpO2, nhiệt độ và độ ẩm. Dữ liệu được truyền lên tầng cloud, hiển thị trên ứng dụng Android và được kết hợp với mô-đun tư vấn sức khỏe sử dụng Firebase AI/Gemini.

Khác với một ứng dụng đo cảm biến đơn giản, hệ thống được thiết kế theo mô hình nhiều tầng có xét đến bảo mật ngay từ quá trình provisioning. Thiết bị ESP32-S3 sử dụng BLE provisioning theo cơ chế Security 1 của Espressif, có Proof of Possession để hạn chế thiết bị hoặc người dùng không hợp lệ tham gia cấu hình mạng. Ở tầng backend, hệ thống Golang Gin triển khai đăng ký, đăng nhập, băm mật khẩu bằng BCrypt, phát hành JWT và mô phỏng luồng OAuth 2.0 Device Authorization Grant cho thiết bị IoT không có giao diện nhập liệu. Redis được dùng làm kho phiên tạm thời có TTL 300 giây, MongoDB dùng lưu dữ liệu người dùng, thiết bị và lịch sử telemetry theo Hourly Bucket Pattern, còn EMQX đóng vai trò MQTT broker cho luồng dữ liệu sinh hiệu.

Báo cáo trình bày toàn bộ quá trình thiết kế kiến trúc, cơ sở lý thuyết của các giao thức sử dụng, hiện thực hóa trên firmware, ứng dụng Android và backend, đồng thời phân tích các rủi ro bảo mật như nghe lén provisioning, chiếm quyền ghép cặp thiết bị, giả mạo telemetry, lộ dữ liệu y tế và tấn công spam vào API.

\tableofcontents
\listoffigures
\listoftables
\newpage
\pagenumbering{arabic}
\pagestyle{fancy}
\fancyhf{}
\fancyhead[R]{\leftmark}
\fancyfoot[R]{\thepage}

\chapter{TỔNG QUAN ĐỀ TÀI}
\section{Giới thiệu đề tài}
Các hệ thống IoT y tế hiện đại không chỉ dừng lại ở việc đọc cảm biến và gửi dữ liệu lên điện thoại. Khi dữ liệu liên quan trực tiếp tới sức khỏe cá nhân, kiến trúc hệ thống phải giải quyết đồng thời nhiều yêu cầu: thiết bị phải dễ cấu hình đối với người dùng phổ thông, dữ liệu phải được cập nhật gần thời gian thực, lịch sử đo phải được lưu trữ theo cách tối ưu, và đặc biệt là quá trình xác thực thiết bị phải đủ an toàn để tránh kẻ tấn công chiếm quyền ghép cặp hoặc giả mạo dữ liệu sinh hiệu.

Trong thực tế, một thiết bị ESP32 đặt trong môi trường gia đình thường không có màn hình, bàn phím hoặc giao diện nhập liệu. Nếu hard-code SSID, mật khẩu Wi-Fi và khóa API vào firmware, hệ thống trở nên khó triển khai hàng loạt và nguy hiểm khi thiết bị bị trích xuất firmware. Nếu bỏ qua bước xác thực thiết bị, một attacker có thể tự tạo MAC giả và gửi telemetry độc hại lên backend. Nếu ghi từng điểm dữ liệu sinh hiệu thành một document riêng, cơ sở dữ liệu sẽ nhanh chóng phình to khi thiết bị gửi dữ liệu liên tục. Các vấn đề này khiến bài toán IoT y tế trở thành một bài toán kiến trúc và bảo mật đúng nghĩa.

\section{Phát biểu vấn đề}
Khi triển khai các ứng dụng IoT trong lĩnh vực y tế, kỹ sư hệ thống phải đối mặt với một nghịch lý lớn: thiết bị biên có tài nguyên tính toán và năng lượng hạn chế, trong khi dữ liệu y tế lại đòi hỏi mức bảo mật đa tầng và nghiêm ngặt hơn nhiều so với dữ liệu cảm biến thông thường. Từ bản thiết kế đang xây dựng, có ba nhóm vấn đề cốt lõi cần được giải quyết.

\subsection{Bề mặt tấn công trong quy trình kích hoạt thiết bị}
Nhiều hệ thống IoT thương mại hoặc đồ án thử nghiệm thường đơn giản hóa quá trình provisioning bằng cách hard-code SSID, mật khẩu Wi-Fi, token Firebase hoặc chỉ dùng địa chỉ MAC công khai để định danh thiết bị. Cách làm này mở ra hai nguy cơ lớn. Thứ nhất là nghe lén thông tin cấu hình nếu quá trình truyền dữ liệu Wi-Fi không được mã hóa. Thứ hai là giả mạo thiết bị, vì địa chỉ MAC có thể bị thu thập qua quét sóng vô tuyến và bị dùng để gửi dữ liệu sinh hiệu giả lên server. Vì vậy, hệ thống cần một quy trình kích hoạt có secure session, Proof of Possession và cơ chế ủy quyền thiết bị phù hợp với thiết bị không có màn hình, tương tự hướng tiếp cận RFC 8628.

\subsection{Nguy cơ lộ lọt dữ liệu y tế trên đường truyền}
Dữ liệu nhịp tim, SpO2 và lịch sử tư vấn sức khỏe có thể được xem là thông tin cá nhân nhạy cảm. Nếu dữ liệu được truyền qua HTTP hoặc MQTT port 1883 không mã hóa, attacker trong cùng mạng có thể thực hiện Man-in-the-Middle để đọc, ghi lại hoặc nguy hiểm hơn là sửa đổi chỉ số sinh hiệu trước khi dữ liệu tới backend. Trong bối cảnh y tế, vi phạm tính toàn vẹn của dữ liệu có thể dẫn tới cảnh báo sai hoặc bỏ sót tình huống nguy hiểm.

\subsection{Cạn kiệt tài nguyên và tấn công DDoS/Spam API}
Thiết bị IoT có thể gửi dữ liệu theo từng giây, tạo áp lực lớn lên database IOPS và dung lượng lưu trữ. Nếu mỗi điểm đo được ghi thành một document riêng, cơ sở dữ liệu nhanh chóng phình to, index lớn và truy vấn lịch sử chậm. Mặt khác, các API công khai phục vụ kích hoạt thiết bị có thể bị spam để tạo hàng nghìn phiên giả, làm cạn RAM Redis hoặc CPU backend. Do đó, thiết kế cần TTL cho phiên tạm, bucket hóa telemetry và cơ chế lọc/rate limit ở tầng cloud.
\section{Mục tiêu nghiên cứu và phạm vi hiện thực}
\textbf{Thiết kế tầng thiết bị nhúng:} ESP32-S3 đọc dữ liệu từ MAX30102 và DHT22, quản lý phiên đo 15 giây, phát hiện trạng thái đặt ngón tay, đồng bộ dữ liệu giữa các task FreeRTOS và upload dữ liệu lên cloud.

\textbf{Thiết kế tầng provisioning và ứng dụng di động:} Android app quét BLE, thiết lập secure session với ESP32, gửi thông tin Wi-Fi, lắng nghe dữ liệu realtime từ Firebase RTDB, đồng thời cung cấp giao diện tư vấn sức khỏe bằng AI.

\textbf{Thiết kế tầng backend và dữ liệu:} Backend Golang Gin cung cấp API đăng ký/đăng nhập, luồng cấp quyền thiết bị theo Device Authorization Grant, MQTT worker xử lý telemetry và lưu dữ liệu vào MongoDB theo bucket.

\textbf{Phân tích bảo mật:} Nhận diện các điểm tấn công trên từng tầng, bao gồm BLE provisioning, API authentication, Device Flow, MQTT broker, Firebase rule và lưu trữ dữ liệu y tế.

\section{Cấu trúc báo cáo}
Báo cáo được tổ chức thành bảy chương. Chương 1 giới thiệu bối cảnh, mục tiêu, phạm vi và phát biểu vấn đề của đề tài. Chương 2 trình bày kiến trúc tổng thể của hệ thống, bao gồm thiết bị ESP32-S3, ứng dụng Android, backend Go và các thành phần hạ tầng dữ liệu. Chương 3 trình bày cơ sở lý thuyết và thiết kế dữ liệu, tập trung vào đo sinh hiệu, BLE provisioning, MQTT, Device Authorization Grant, Redis TTL và MongoDB Bucket Pattern. Chương 4 đi sâu vào quá trình triển khai hệ thống từ mã nguồn firmware, Android app, backend và Docker Compose. Chương 5 mô tả quy trình vận hành và các luồng dữ liệu chính, bao gồm provisioning, ủy quyền thiết bị, telemetry và tư vấn AI. Chương 6 thảo luận kết quả kiểm thử, giao diện giám sát và phân tích bảo mật hệ thống. Chương 7 tổng kết kết quả đạt được, nêu hạn chế và đề xuất hướng phát triển.

\chapter{KIẾN TRÚC HỆ THỐNG}
\section{Tổng quan về kiến trúc hệ thống}
Kiến trúc hệ thống được xây dựng theo mô hình phân tầng công nghệ kết hợp với cơ chế định tuyến hướng sự kiện. Mô hình này đảm bảo tính tách biệt (Separation of Concerns), tối ưu hóa luồng dữ liệu thời gian thực và thực thi các lớp bảo mật đa tầng. Toàn bộ hệ thống vận hành dựa trên sự tương tác chặt chẽ của 4 thành phần cốt lõi:

\begin{itemize}
    \item \textbf{ESP32 (Edge Device):} Đóng vai trò là thành phần phần cứng tại biên (Edge Layer), giao tiếp trực tiếp với các mô-đun cảm biến y tế để thu thập chỉ số sinh hiệu bao gồm nhịp tim và nồng độ oxy trong máu ($SpO_2$). Thiết bị chịu trách nhiệm xử lý tín hiệu số, thực hiện thuật toán tiền xử lý dữ liệu thô và nén dữ liệu thành chuỗi ký tự tối giản nhằm tối ưu hóa băng thông truyền thông. Về mặt an ninh, ESP32 tích hợp các công cụ mã hóa phần cứng để thực hiện cơ chế bắt tay trao đổi khóa ECDH, mã hóa đối xứng AES-128 trong pha thiết lập cục bộ và khởi tạo kênh truyền bảo mật MQTT over TLS trước khi đẩy dữ liệu lên hệ thống Cloud qua kết nối mạng không dây Wi-Fi.
    
    \item \textbf{MQTT Broker (EMQX):} Đóng vai trò là trung tâm điều phối, phân phối và luân chuyển thông điệp ở tầng mạng (Network Layer) theo mô hình kiến trúc Publish/Subscribe. Để đảm bảo tính toàn vẹn và bí mật của dữ liệu y tế, EMQX được cấu hình lớp bảo mật Transport Layer Security (TLS) nhằm mã hóa toàn diện đường truyền từ biên lên Cloud. Đồng thời, Broker áp dụng cơ chế kiểm soát truy cập ACL (Access Control List) chặt chẽ để phân quyền định tuyến riêng biệt cho từng thiết bị theo địa chỉ MAC phần cứng. Kiến trúc phân tán của EMQX giúp hệ thống đạt khả năng chịu tải cao, giảm thiểu tối đa độ trễ truyền dẫn và tăng cường năng lực phòng chống các cuộc tấn công từ chối dịch vụ (DDoS) trực diện vào tầng mạng.
    
    \item \textbf{Backend Go (Core Service):} Được định vị là trung tâm xử lý logic (Core Business Logic) của toàn bộ hệ thống Cloud, phát triển trên nền tảng ngôn ngữ Golang kết hợp với Gin Framework nhờ ưu thế vượt trội về hiệu năng xử lý đa luồng (Concurrency). Go Backend đóng vai trò là một Subscriber bất đồng bộ kết nối trực tiếp vào EMQX để đón nhận và phân tách các luồng dữ liệu (Ingress Streams) truyền về từ thiết bị biên. Thành phần này chịu trách nhiệm quản lý máy trạng thái (State Machine) của giao thức xác thực thiết bị không đầu theo tiêu chuẩn RFC 8628, điều phối phiên làm việc thông qua bộ nhớ đệm Redis và tổ chức phân phối dữ liệu sinh hiệu vào cơ sở dữ liệu lưu trữ lịch sử dài hạn.
    
    \item \textbf{Firebase Realtime Database (RTDB):} Đóng vai trò là cơ sở dữ liệu đồng bộ trạng thái thời gian thực và là cầu nối trực tiếp đến ứng dụng phía người dùng (App Client). Sau khi dữ liệu sinh hiệu được bóc tách và xử lý tại Backend Go, hệ thống sẽ đẩy trực tiếp dữ liệu này lên Firebase RTDB thông qua giao thức kết nối liên tục WebSockets. Cơ chế hướng sự kiện này giúp ứng dụng Android phía Client tự động cập nhật và hiển thị chỉ số tức thì với độ trễ phản hồi cực thấp (dưới $100\text{ ms}$) mà không cần thực hiện các kỹ thuật kéo dữ liệu định kỳ (Polling/Pull). Giải pháp này giảm thiểu tối đa tài nguyên xử lý của CPU, tiết kiệm dung lượng pin và tối ưu hóa băng thông mạng cho thiết bị di động của người dùng.
\end{itemize}

\begin{figure}[H]
    \centering
    % Sử dụng dấu ngoặc kép bọc quanh tên file để trình biên dịch xử lý được khoảng trắng
    \includegraphics[width=0.95\textwidth]{"ESP32 Data Processing-2026-07-10-072915"}
    \caption{Sơ đồ luồng xử lý và kiến trúc tương tác giữa 4 thành phần cốt lõi.}
    \label{fig:esp32_data_processing}
\end{figure}

\section{Luồng vận hành của hệ thống}
Để đảm bảo tính toàn vẹn dữ liệu, an ninh bảo mật và khả năng đáp ứng thời gian thực, các thành phần trong kiến trúc tương tác chặt chẽ với nhau thông qua hai luồng nghiệp vụ cốt lõi: Luồng thiết lập \& ủy quyền thiết bị lần đầu (Provisioning) và Luồng nhận \& ghi dữ liệu sinh hiệu hàng ngày (Telemetry).

\subsection{Luồng Thiết lập \& Ủy quyền Thiết bị Lần đầu (Provisioning - RFC 8628)}
Quá trình khởi tạo và gắn định danh thiết bị vào tài khoản người dùng được chia làm hai giai đoạn độc lập nhằm tối ưu hóa tính tiện dụng và an toàn bảo mật:

\subsubsection{Giai đoạn cấu hình mạng thông qua BLE}
\begin{enumerate}
    \item Khi thiết bị ESP32-S3 khởi động lần đầu hoặc được kích hoạt chế độ cấu hình, hệ thống sẽ phát sóng Bluetooth Low Energy (BLE) với tên định danh mang tiền tố \texttt{Prov\_}.
    \item Người dùng sử dụng Ứng dụng Android thực hiện quét mã QR đính kèm trên thiết bị để thu thập thông tin định danh và địa chỉ MAC. 
    \text{Ứng dụng} thiết lập một phiên kết nối bảo mật với thiết bị biên thông qua giao thức bắt tay trao đổi khóa công khai ECDH (Elliptic Curve Diffie-Hellman).
    \item Sau khi thống nhất khóa đối xứng tạm thời, ứng dụng di động tiến hành mã hóa thông tin cấu hình mạng Wi-Fi (gồm SSID và mật khẩu) bằng thuật toán AES-128 rồi truyền sang ESP32-S3. Thiết bị nhận dữ liệu, giải mã, thực hiện kết nối vào mạng Wi-Fi cục bộ và chủ động ngắt hoàn toàn sóng BLE để tiết kiệm năng lượng và triệt tiêu nguy cơ bị tấn công qua kênh vật lý.
\end{enumerate}

\subsubsection{Giai đoạn xác thực và cấp quyền thiết bị theo chuẩn RFC 8628}
\begin{enumerate}
    \item Ngay sau khi thiết lập kết nối Internet thành công, ESP32-S3 sử dụng mã định danh nặc danh gửi một yêu cầu đăng ký thiết bị (\textit{Device Authorization Request}) thông qua giao thức HTTPS REST API tới Go Backend. Yêu cầu này bắt buộc đính kèm mã định danh dòng máy (\texttt{client\_id}) và địa chỉ phần cứng (\texttt{mac\_address}) để Backend thực hiện đối sánh danh mục thiết bị hợp lệ đã xuất xưởng.
    
    \item Go Backend tiếp nhận yêu cầu, tiến hành xác thực thiết bị và sinh ra các tham số theo đúng đặc tả RFC 8628 bao gồm: \texttt{device\_code} (mã định danh phiên bí mật dùng cho thiết bị), \texttt{user\_code} (mã kích hoạt ngắn thân thiện với người dùng, định dạng \texttt{XX-XX-XX}), và \texttt{verification\_uri} (đường dẫn trang kích hoạt). Tập hợp các tham số này cùng trạng thái khởi tạo \texttt{pending} được lưu giữ trong bộ nhớ đệm Redis Cache với thời gian sống nghiêm ngặt $\text{TTL} = 300\text{ s}$ nhằm chống tấn công vét cạn. Đồng thời, Backend đồng bộ \texttt{user\_code} này lên Firebase RTDB dưới một nút phiên làm việc được liên kết bởi địa chỉ MAC của thiết bị.
    
    \item Nhằm tối ưu hóa trải nghiệm người dùng và loại bỏ bước nhập liệu thủ công của chuẩn RFC 8628 gốc, Ứng dụng Android (đã thiết lập kết nối liên tục WebSockets với Firebase RTDB làm kênh truyền phụ - \textit{Out-of-band Channel}) sẽ tự động phát hiện và bắt trọn mã \texttt{user\_code} dựa vào địa chỉ MAC tương ứng đã ghi nhận ở pha BLE. Màn hình ứng dụng lập tức hiển thị giao diện yêu cầu xác nhận kích hoạt thiết bị y tế mà không cần người dùng nhập tay.
    
    \item Khi người dùng nhấn nút phê duyệt trên ứng dụng, ứng dụng gửi một yêu cầu xác nhận kèm Token JWT định danh của chính người dùng lên API phê duyệt của Go Backend tại đường dẫn: \url{/api/v1/auth/user/confirm}. Backend tiến hành giải mã chữ ký JWT để xác định danh tính người dùng, kiểm tra tính hợp lệ của phiên trên Redis thông qua \texttt{user\_code}, sau đó cập nhật trạng thái phiên từ \texttt{pending} sang \texttt{approved} và liên kết phiên này với tài khoản ID của người dùng.
    
    \item Song song với quá trình trên, kể từ sau bước 1, ESP32-S3 liên tục thực hiện cơ chế gửi yêu cầu thăm dò định kỳ (\textit{Device Access Token Polling}) sau mỗi khoảng thời gian $\Delta t = 5\text{ s}$ thông qua API \url{/api/v1/auth/token} bằng cách gửi lên \texttt{device\_code}. Tại lần thăm dò tiếp theo ngay sau khi phiên chuyển sang trạng thái \texttt{approved}, Go Backend sẽ tiến hành sinh ra một Access Token dài hạn kết hợp chữ ký số bảo mật và trả về cho ESP32-S3 trong phản hồi HTTP 200 OK, đồng thời giải phóng phiên tạm trên Redis. Thiết bị biên nhận Access Token, lưu vào bộ nhớ Flash và kết thúc luồng Polling để chuyển sang luồng đẩy dữ liệu sinh hiệu.
\end{enumerate}

\subsection{Luồng Nhận \& Ghi dữ liệu Sinh hiệu Hàng ngày (Telemetry Pipeline)}
Sau khi hoàn tất giai đoạn ủy quyền, luồng truyền nhận và xử lý dữ liệu sinh hiệu thời gian thực được vận hành tuần tự theo mô hình hướng sự kiện:

\begin{enumerate}
    \item \textbf{Thu thập và nén tại biên:} Định kỳ mỗi giây ($f = 1\text{ Hz}$), các mô-đun cảm biến (MAX30102, DHT22) hoàn tất quá trình biến đổi tương tự - số (ADC) và chuyển dữ liệu thô về vi điều khiển ESP32-S3 qua giao tiếp ngoại vi ($I^2C$/One-Wire). Thiết bị thực hiện các thuật toán lọc số, bóc tách chỉ số nhịp tim, $SpO_2$, nhiệt độ và độ ẩm môi trường, sau đó đóng gói thành một chuỗi văn bản nén tối giản nhằm tối ưu hóa payload truyền dẫn theo định dạng: \texttt{"hr,spo2,temp,hum"} (Ví dụ: \texttt{"75,98,32.5,65.0"}).
    
    \item \textbf{Vận chuyển bảo mật tầng mạng:} Chuỗi dữ liệu nén được ESP32-S3 đóng gói vào gói tin mạng của giao thức MQTT. Thiết bị đính kèm mã địa chỉ phần cứng vào trường \texttt{Username} và Access Token (được cấp từ luồng RFC 8628) vào trường \texttt{Password} để thực hiện xác thực, sau đó tiến hành phát tán (\textit{Publish}) lên EMQX Broker thông qua kênh truyền mã hóa toàn diện MQTT over TLS (Cổng bảo mật 8883). Đích đến là một Topic được cô lập nghiêm ngặt bằng cấu hình ACL theo cấu trúc: \texttt{devices/\{MAC\_ADDRESS\}/telemetry}.
    
    \item \textbf{Xử lý bất đồng bộ tại Cloud:} EMQX Broker kiểm tra quyền ACL, chấp thuận thông điệp và ngay lập tức chuyển tiếp (\textit{Forward}) luồng dữ liệu tới Go Backend thông qua một MQTT Worker chạy ẩn (\textit{Background Ingress Worker}) dựa trên cơ chế non-blocking. Go Backend tiếp nhận thông điệp, bóc tách chuỗi văn bản \texttt{"75,98,32.5,65.0"} thành các trường dữ liệu số độc lập và thực hiện phân phối bất đồng bộ song song làm hai nhánh xử lý:
    
    \begin{itemize}
        \item \textbf{Nhánh đồng bộ thời gian thực:} Go Backend tương tác với dịch vụ đám mây Firebase thông qua \textit{Firebase Admin SDK} để ghi trực tiếp dữ liệu sinh hiệu vừa bóc tách vào nút cây lưu trữ tương ứng của Firebase RTDB. Nhờ kết nối WebSockets duy trì sẵn giữa Firebase và Client, hệ thống lập tức cập nhật trạng thái và đẩy (\textit{Push}) dữ liệu xuống ứng dụng Android phía Client với độ trễ phản hồi cực thấp ($t < 100\text{ ms}$), giúp biểu đồ sinh hiệu trên ứng dụng nhảy số mượt mà theo thời gian thực.
        
        \item \textbf{Nhánh lưu trữ lịch sử dài hạn:} Nhằm tối ưu hóa hiệu năng I/O đĩa và tiết kiệm dung lượng lưu trữ, Go Backend không thực hiện ghi riêng lẻ từng bản ghi theo từng giây vào cơ sở dữ liệu. Thay vào đó, hệ thống áp dụng giải pháp thiết kế \textbf{Bucket Pattern}. Dữ liệu sinh hiệu của từng thiết bị được gom cụm lại theo đơn vị giờ và ghi dồn mảng (\textit{Push array element}) vào một tài liệu (\textit{Document}) duy nhất trong cơ sở dữ liệu MongoDB tương ứng với khung giờ hiện hành, giúp giảm thiểu số lượng truy vấn và tăng tốc độ đọc dữ liệu lịch sử sau này.
    \end{itemize}
\end{enumerate}

\section{Thiết kế phần cứng đo sinh hiệu}
Thiết bị đo sử dụng ESP32-S3 DevKitC-1 làm vi điều khiển trung tâm. MAX30102 được kết nối qua bus I2C để thu tín hiệu hồng ngoại và đỏ phục vụ phát hiện nhịp tim và ước lượng SpO2. DHT22 đo nhiệt độ và độ ẩm môi trường, giúp app cung cấp thêm bối cảnh sức khỏe khi người dùng hỏi AI. LED trên chân GPIO 2 được dùng làm chỉ báo trạng thái provisioning.

\begin{table}[H]
\centering
\caption{Bảng thông số phần cứng và chân kết nối chính}
\renewcommand{\arraystretch}{1.25}
\setlength{\tabcolsep}{5pt}
\begin{tabularx}{\textwidth}{p{3cm}p{5cm}X}
\toprule
\textbf{Thành phần} & \textbf{Thông số trong code} & \textbf{Vai trò} \\
\midrule
ESP32-S3 & Board \code{esp32-s3-devkitc-1} & Vi điều khiển trung tâm, hỗ trợ Wi-Fi, BLE, dual-core và FreeRTOS. \\
MAX30102/\newline MAX30105\newline library & I2C SDA = 8, SCL = 9 & Đọc tín hiệu IR/Red để phát hiện nhịp tim và ước lượng SpO2. \\
DHT22 & DATA = GPIO 4 & Đo nhiệt độ và độ ẩm môi trường. \\
LED trạng thái & GPIO 2 & Báo hiệu chế độ provisioning qua BLE. \\
Firebase client & \texttt{Firebase\_ESP\_Client} & Upload dữ liệu realtime và lịch sử lên Firebase RTDB. \\
\bottomrule
\end{tabularx}
\end{table}

Thiết bị không đo liên tục như một máy y tế chuyên dụng mà tổ chức dữ liệu theo phiên đo. Khi cảm biến phát hiện ngón tay với ngưỡng IR lớn hơn 50000, firmware chuyển từ trạng thái \texttt{STATE\_IDLE} sang \texttt{STATE\_MEASURING}. Trong 15 giây, các giá trị BPM và SpO2 được cộng dồn để lấy trung bình. Khi hết phiên đo, thiết bị chuyển sang \texttt{STATE\_COMPLETED} và gửi dữ liệu lên Firebase theo hai nhánh:

\begin{itemize}
    \item Kết quả mới nhất được ghi tại \texttt{/devices/\{mac\}/latest}.
    \item Bản ghi lịch sử được lưu tại \texttt{/devices/\{mac\}/history/\{timestamp\}}.
\end{itemize}

\section{Kiến trúc ứng dụng Android}
Ứng dụng Android được tổ chức theo Clean Architecture kết hợp MVVM. Tầng giao diện sử dụng Jetpack Compose và Material 3, tầng ViewModel quản lý state bằng StateFlow, tầng domain định nghĩa các model và use case, còn tầng data giao tiếp với Room, Firebase AI, Firebase RTDB, BLE service và Location service.

\textbf{BLE provisioning service:} Lớp \texttt{BleIoTService} sử dụng Espressif Provisioning SDK để quét thiết bị BLE có tên bắt đầu bằng \texttt{Prov\_}. Sau khi phát hiện thiết bị, app tạo \texttt{ESPDevice} với transport BLE và Security 1, gán Proof of Possession mặc định \texttt{12345678}, kết nối tới service UUID và gửi SSID/mật khẩu Wi-Fi qua secure session.

\textbf{Realtime IoT ViewModel:} Lớp \texttt{IoTViewModel} lắng nghe Firebase RTDB ở đường dẫn \texttt{devices/\{deviceId\}/latest}. Khi BLE service phát hiện MAC mới, ViewModel gỡ listener cũ và chuyển sang đường dẫn của thiết bị mới. Dữ liệu được ánh xạ vào model \texttt{IoTData}, gồm nhịp tim, SpO2, nhiệt độ, độ ẩm, trạng thái đặt ngón tay, trạng thái phiên đo và lịch sử nhịp tim ngắn hạn.

\textbf{AI medical assistant:} Lớp \texttt{ChatManager} giao tiếp với Firebase AI/Gemini. Ứng dụng lưu lịch sử hội thoại, nén ngữ cảnh khi dài, trích xuất symptom cache và kết hợp thông tin địa điểm y tế gần người dùng để đưa ra phản hồi phù hợp. Đây là phần mở rộng giúp hệ thống không chỉ hiển thị số đo mà còn hỗ trợ người dùng diễn giải triệu chứng ở mức tham khảo.

\section{Kiến trúc backend và hạ tầng dữ liệu}
Backend Golang Gin được triển khai theo các module rõ ràng. \texttt{main.go} khởi tạo MongoDB, Redis, MQTT worker và router HTTP. \texttt{auth\_handler.go} xử lý đăng ký/đăng nhập. \texttt{device\_handler.go} xử lý Device Authorization Flow. \texttt{mqtt\_worker.go} subscribe topic telemetry và ghi dữ liệu lịch sử. \texttt{database.go} định nghĩa interface \texttt{DatabaseService} để test có thể mock database độc lập.

\begin{table}[H]
\centering
\caption{Vai trò của các thành phần backend}
\begin{tabular}{p{4cm}p{10cm}}
\toprule
\textbf{Thành phần} & \textbf{Vai trò trong hệ thống} \\
\midrule
Gin HTTP API & Cung cấp endpoint đăng ký, đăng nhập, authorize, confirm và token. \\
Redis & Lưu phiên cấp quyền thiết bị tạm thời với TTL 300 giây. \\
MongoDB & Lưu user, device đã ghép cặp và lịch sử telemetry theo bucket. \\
EMQX & MQTT broker nhận telemetry từ thiết bị theo topic \texttt{devices/+/telemetry}. \\
Firebase RTDB & Kênh realtime cho app Android; backend hiện mới mô phỏng ghi RTDB qua log. \\
Docker Compose & Đóng gói backend, MongoDB, Redis và EMQX để chạy local đồng bộ. \\
\bottomrule
\end{tabular}
\end{table}

\chapter{CƠ SỞ LÝ THUYẾT}
\section{Xử lý tín hiệu nhịp tim và SpO2 trên thiết bị nhúng}
Cảm biến MAX30102 hoạt động dựa trên nguyên lý quang học PPG (Photoplethysmography). LED hồng ngoại và LED đỏ chiếu qua mô ngón tay, photodiode đo lượng ánh sáng phản xạ. Khi tim co bóp, thể tích máu trong mao mạch thay đổi làm tín hiệu hấp thụ ánh sáng biến thiên tuần hoàn. Firmware sử dụng hàm \texttt{checkForBeat()} từ thư viện SparkFun để phát hiện đỉnh nhịp và tính BPM theo khoảng thời gian giữa hai nhịp liên tiếp:
\begin{equation}
BPM = \frac{60}{\Delta t_{beat}}
\end{equation}
Trong đó $\Delta t_{beat}$ là khoảng thời gian giữa hai lần phát hiện nhịp, tính bằng giây. Để giảm nhiễu, firmware chỉ nhận các giá trị BPM trong khoảng 40 đến 180 và lấy trung bình trượt với kích thước \texttt{RATE\_SIZE = 4}.

Đối với SpO2, code hiện thực công thức xấp xỉ đơn giản dựa trên tỷ lệ tín hiệu đỏ và hồng ngoại:
\begin{equation}
R = \frac{Red}{IR}, \qquad SpO2 \approx 110 - 25R
\end{equation}
Giá trị SpO2 được giới hạn trong khoảng 80 đến 100 để tránh hiển thị các kết quả phi thực tế do nhiễu cảm biến hoặc đặt tay sai vị trí. Cách tính này phù hợp với nguyên mẫu học phần nhưng chưa thể thay thế thuật toán hiệu chuẩn y tế của thiết bị thương mại.

\section{BLE provisioning và Proof of Possession}
Provisioning là quá trình đưa thông tin Wi-Fi từ app sang thiết bị. Nếu truyền SSID và mật khẩu dạng plaintext qua BLE, attacker ở gần có thể nghe lén. Espressif Security 1 dùng cơ chế bắt tay bảo mật dựa trên trao đổi khóa, sau đó mã hóa dữ liệu cấu hình. Proof of Possession (PoP) là bí mật ngắn mà cả app và thiết bị cùng biết, nhằm chứng minh người dùng đang cấu hình đúng thiết bị vật lý.

Trong project, PoP đang được đặt là \texttt{12345678}. Về mặt kiến trúc, biến này nên được thay bằng mã riêng theo từng thiết bị, in trên tem hoặc mã QR. Nếu toàn bộ thiết bị dùng chung một PoP, attacker biết PoP của một thiết bị có thể thử cấu hình các thiết bị khác cùng dòng.

\section{OAuth 2.0 Device Authorization Grant cho IoT}
Thiết bị IoT thường không có trình duyệt hoặc bàn phím, vì vậy luồng đăng nhập OAuth truyền thống không phù hợp. OAuth 2.0 Device Authorization Grant cho phép thiết bị khởi tạo yêu cầu ủy quyền, nhận \texttt{device\_code} và \texttt{user\_code}; người dùng dùng app hoặc trình duyệt để xác nhận; thiết bị polling backend cho tới khi được cấp token.

Trong backend của đề tài, luồng này được triển khai theo ba endpoint chính: authorize để sinh mã và lưu Redis, confirm để app duyệt phiên bằng JWT và chữ ký HMAC, token để ESP32 polling và nhận access token sau khi phiên được duyệt.

\section{MQTT, topic ACL và mô hình telemetry}
MQTT phù hợp với IoT vì header nhẹ, hỗ trợ publish/subscribe và cho phép backend tách rời khỏi số lượng thiết bị. Trong hệ thống, topic telemetry được thiết kế theo MAC của từng thiết bị: \texttt{devices/\{MAC\}/telemetry}

Payload hiện tại là chuỗi ngắn theo dạng \texttt{bpm,spo2}. Ví dụ, \texttt{75,98} nghĩa là BPM = 75 và SpO2 = 98. MQTT worker kiểm tra format, parse số nguyên và đánh giá trạng thái \texttt{Normal} nếu BPM nằm trong khoảng 60--100 và SpO2 không thấp hơn 95. Trường hợp còn lại được đánh dấu là \texttt{Warning}.

Khi triển khai thực tế, broker cần bật TLS và ACL để mỗi device token chỉ được publish lên topic của chính MAC tương ứng.

\section{Redis TTL và MongoDB Hourly Bucket Pattern}
Redis được dùng cho dữ liệu tạm thời vì tốc độ cao và hỗ trợ TTL tự động. Mỗi phiên Device Flow được lưu dưới key \texttt{device\_flow:\{user\_code\}}. TTL 300 giây có ba ý nghĩa bảo mật: giới hạn thời gian attacker brute-force user code, tránh replay phiên cũ và tự dọn bộ nhớ khi thiết bị bỏ dở quá trình ghép cặp.

MongoDB lưu telemetry theo Hourly Bucket Pattern. Thay vì tạo một document cho mỗi điểm đo, backend gom các điểm đo cùng thiết bị, cùng ngày và cùng giờ vào một document có khóa:
\begin{equation}
BucketID = MAC \Vert Date \Vert Hour
\end{equation}
Mỗi bucket chứa mảng \texttt{data\_points}. Mô hình này giảm số lần tạo document mới, tối ưu truy vấn theo thời gian và phù hợp với dữ liệu sinh hiệu gửi liên tục.

\section{Đặc tả Redis, MongoDB và API theo bản thiết kế}
\subsection{Redis Cache cho phiên xác thực tạm thời}
Trong thiết kế hệ thống, Redis được dùng như lớp in-memory cache tốc độ cao cho luồng RFC 8628. Key có dạng logic \texttt{auth:device\_code:\{device\_code\}} hoặc trong code hiện tại là \texttt{device\_flow:\{user\_code\}}. Value là JSON chứa \texttt{user\_code}, định danh thiết bị như \texttt{device\_id} hoặc MAC, \texttt{session\_id}, trạng thái \texttt{pending}/\texttt{approved} và token sau khi được duyệt.. TTL 300 giây là lựa chọn hợp lý vì vừa giảm rủi ro brute-force, vừa tránh Redis bị tràn RAM khi có nhiều phiên giả hoặc thiết bị kích hoạt đồng thời.

\begin{lstlisting}
{
  "user_code": "X7-R9-B2",
  "device_id": "ESP32_S3_MAC_A1B2C3",
  "session_id": "sess_9876543210",
  "status": "pending",
  "access_token": ""
}
\end{lstlisting}

\subsection{MongoDB Bucket Pattern cho lịch sử sinh hiệu}
Nếu ghi từng mẫu đo thành một document riêng, một thiết bị gửi 1 mẫu/giây sẽ tạo 3.600 document mỗi giờ. Với nhiều thiết bị, index tăng nhanh và truy vấn biểu đồ lịch sử trở nên tốn CPU/I/O. Bucket Pattern gom toàn bộ sample cùng thiết bị, cùng ngày và cùng giờ vào một document duy nhất. Nhờ vậy, mỗi thiết bị chỉ tạo 24 document/ngày thay vì 86.400 document/ngày nếu ghi theo giây.

\begin{table}[H]
\centering
\caption{So sánh ghi theo từng giây và Bucket Pattern}
\begin{tabular}{p{4cm}p{5cm}p{5cm}}
\toprule
\textbf{Tiêu chí} & \textbf{Ghi từng giây} & \textbf{Bucket Pattern theo giờ} \\
\midrule
Số document & 3.600 document/giờ/thiết bị & 1 document/giờ/thiết bị \\
Index & Lớn nhanh, dễ vượt RAM cache & Nhỏ gọn, truy vấn theo giờ hiệu quả \\
Truy vấn biểu đồ & Phải quét nhiều bản ghi nhỏ & Đọc một bucket là đủ dữ liệu 1 giờ \\
Dung lượng & Lặp metadata ở mỗi document & Metadata chỉ lưu một lần trong bucket \\
\bottomrule
\end{tabular}
\end{table}

\subsection{API Device Flow và MQTT ACL}
Luồng Device Flow của hệ thống gồm ba API logic chính:

\begin{itemize}
    \item Yêu cầu cấp mã thiết bị: \url{/api/v1/oauth/device/authorize}.
    \item Polling token: \url{/api/v1/oauth/token}.
    \item Người dùng xác nhận trên ứng dụng: \url{/api/v1/oauth/device/confirm}.
\end{itemize}

Các endpoint này tạo thành luồng xác thực thiết bị hoàn chỉnh: thiết bị xin mã, backend lưu phiên tạm thời, ứng dụng người dùng xác nhận, sau đó thiết bị nhận token để giao tiếp với hệ thống.

Đối với MQTT, topic chuẩn là \texttt{devices/\{MAC\}/telemetry}. EMQX cần ACL theo nguyên tắc: thiết bị chỉ được publish vào topic của chính nó, không được subscribe dữ liệu của thiết bị khác; backend worker được subscribe \texttt{devices/+/telemetry}. Đây là lớp bảo vệ quan trọng để tránh việc một thiết bị bị chiếm quyền đọc toàn bộ dữ liệu sinh hiệu trong hệ thống.
\chapter{TRIỂN KHAI HỆ THỐNG}
\section{Hiện thực firmware ESP32-S3}
Firmware trong \texttt{medical-ai-chatbot-firmware/src/main.cpp} được tổ chức theo các nhóm chức năng rõ ràng: cấu hình chân kết nối, khai báo cảm biến, biến tính nhịp tim, biến DHT22, trạng thái phiên đo, dữ liệu chia sẻ giữa hai core, BLE provisioning và task upload Firebase.

\begin{table}[H]
\centering
\caption{Máy trạng thái phiên đo trên ESP32-S3}
\begin{tabular}{p{4cm}p{10cm}}
\toprule
\textbf{Trạng thái} & \textbf{Ý nghĩa và hành vi} \\
\midrule
\texttt{STATE\_IDLE} & Chưa phát hiện ngón tay. Thiết bị vẫn gửi gói idle định kỳ để app biết thiết bị còn hoạt động. \\
\texttt{STATE\_MEASURING} & Đã phát hiện ngón tay. Firmware bắt đầu phiên 15 giây, tích lũy BPM và SpO2 để lấy trung bình. \\
\texttt{STATE\_COMPLETED} & Phiên đo kết thúc. Thiết bị gửi giá trị cuối cùng lên latest và ghi thêm một bản lịch sử theo timestamp. \\
\bottomrule
\end{tabular}
\end{table}

\begin{algorithm}[H]
\caption{Thuật toán xử lý phiên đo trên ESP32-S3}
\KwIn{Giá trị IR, Red từ MAX30102; nhiệt độ và độ ẩm từ DHT22}
\KwOut{Bản ghi realtime gửi lên Firebase RTDB}
\While{thiết bị đang chạy}{
    Đọc \texttt{irValue} và \texttt{redValue}\;
    \eIf{\texttt{irValue >= 50000}}{
        Nếu đang \texttt{IDLE}, chuyển sang \texttt{MEASURING} và reset bộ tích lũy\;
        Phát hiện nhịp tim, tính BPM và cập nhật trung bình trượt\;
        Cứ mỗi 2 giây đọc DHT22 và ước lượng SpO2\;
        Nếu đủ 15 giây, chuyển sang \texttt{COMPLETED} và ghi kết quả cuối\;
    }{
        Nếu trước đó đang đo, chuyển về \texttt{IDLE} và gửi trạng thái không có ngón tay\;
    }
    Task Firebase đọc dữ liệu chia sẻ qua mutex và upload bất đồng bộ\;
}
\end{algorithm}

Việc dùng \texttt{xSemaphoreTake()} khi đọc/ghi \texttt{sharedData} là điểm quan trọng trong kiến trúc firmware. ESP32-S3 có hai core, một luồng xử lý cảm biến và một task upload cloud có thể truy cập cùng dữ liệu. Nếu không dùng mutex, app có thể nhận gói dữ liệu không nhất quán, ví dụ BPM của phiên mới nhưng timestamp của phiên cũ.

\section{Hiện thực ứng dụng Android}
Trong Android app, \texttt{BleIoTService} là lớp cầu nối giữa người dùng và thiết bị. Service lấy Bluetooth manager, Wi-Fi manager, quét BLE, nhận sự kiện EventBus từ Espressif SDK và cập nhật trạng thái qua \texttt{MutableStateFlow}. Cách tổ chức này giúp UI Compose chỉ cần observe state thay vì tự xử lý callback Bluetooth phức tạp.

\texttt{IoTViewModel} đóng vai trò như một dashboard dữ liệu nhỏ cho thiết bị sức khỏe. ViewModel lưu trạng thái giao diện bằng \texttt{StateFlow} với kiểu dữ liệu \texttt{IoTData}. Thành phần này cũng tạo listener tới Firebase RTDB để nhận dữ liệu realtime và lọc theo timestamp, nhờ đó tránh việc UI bị cập nhật ngược thời gian. Lịch sử nhịp tim ngắn hạn được giữ tối đa 20 điểm, đủ để vẽ xu hướng nhanh trên màn hình mà không gây nặng bộ nhớ.

Phần AI sử dụng \texttt{ChatManager} để quản lý lịch sử hội thoại với Gemini. Khi hội thoại dài, hàm \texttt{shouldCompress()} kiểm tra số lượng message và tổng ký tự để quyết định nén ngữ cảnh. Ngoài prompt người dùng, request gửi lên model có thể kèm tóm tắt bệnh sử, symptom cache và danh sách cơ sở y tế gần vị trí người dùng. Thiết kế này biến app từ một màn hình đo số đơn thuần thành trợ lý sức khỏe có ngữ cảnh.

\section{Hiện thực backend xác thực và Device Flow}
Backend đăng ký người dùng bằng số điện thoại và mật khẩu. Khi đăng ký, handler kiểm tra số điện thoại đã tồn tại chưa, sau đó băm mật khẩu bằng BCrypt trước khi lưu MongoDB. Khi đăng nhập, backend so sánh mật khẩu bằng \begin{center}
\url{bcrypt.CompareHashAndPassword()}
\end{center}
 và phát hành JWT chứa \texttt{uid\_user}. JWT có thời hạn 30 ngày, phù hợp với trải nghiệm mobile app nhưng khi production nên bổ sung refresh token và revoke token.

Device Flow trong \texttt{device\_handler.go} có một điểm bảo mật quan trọng là kiểm tra chữ ký PoP:
\begin{equation}
Signature = HMAC\_SHA256(userCode : macAddress : sessionId, secret)
\end{equation}
Backend chỉ duyệt phiên nếu chữ ký nhận được trùng với chữ ký tính lại, MAC/session trùng với phiên trong Redis và JWT người dùng hợp lệ. Sau khi thiết bị polling thành công, backend tạo access token dạng \texttt{dev\_token\_...}, lưu vào collection \texttt{devices} và xóa phiên tạm khỏi Redis.

\section{Hiện thực MQTT worker và lưu lịch sử}
File \texttt{mqtt\_worker.go} hiện thực worker nhận telemetry từ MQTT broker. Worker tạo MQTT client cho backend, bật clean session và auto reconnect để có thể tự kết nối lại khi broker gián đoạn. Sau khi kết nối thành công, worker subscribe topic telemetry của các thiết bị theo mẫu \url{devices/+/telemetry}. Mỗi message nhận được được chuyển sang một goroutine riêng để xử lý, nhờ đó callback MQTT không bị chặn bởi các thao tác parse dữ liệu, đánh giá trạng thái hoặc ghi cơ sở dữ liệu.

\begin{algorithm}[H]
\caption{Thuật toán xử lý telemetry trong MQTT worker}
\KwIn{Topic \texttt{devices/\{mac\}/telemetry}, payload \texttt{bpm,spo2}}
\KwOut{Bucket lịch sử MongoDB và cập nhật realtime Firebase mô phỏng}
Tách MAC từ topic\;
Tách payload theo dấu phẩy\;
\If{payload không có đúng hai trường}{ghi log lỗi và bỏ qua}
Parse BPM và SpO2 sang số nguyên\;
\eIf{BPM < 60 hoặc BPM > 100 hoặc SpO2 < 95}{status = \texttt{Warning}}{status = \texttt{Normal}}
Tạo \texttt{TelemetryDataPoint} với timestamp hiện tại\;
Tính \texttt{date} và \texttt{hour}\;
Upsert document \texttt{telemetry\_history} với \texttt{\$push} vào \texttt{data\_points}\;
Gọi hàm mô phỏng cập nhật Firebase RTDB\;
\end{algorithm}

\section{Triển khai bằng Docker Compose}
Hạ tầng backend được đóng gói bằng Docker Compose gồm bốn service: app, mongo, redis và emqx. Cấu hình này giúp nhóm khởi chạy toàn bộ môi trường local bằng một lệnh duy nhất, đồng thời phản ánh cách tách service trong triển khai cloud thật.
\begin{lstlisting}[language=bash]
docker compose up -d --build
\end{lstlisting}
Các biến môi trường \texttt{MONGODB\_URI}, \texttt{REDIS\_ADDR} và \texttt{MQTT\_BROKER\_URI} giúp backend không phụ thuộc vào địa chỉ localhost khi chạy trong container. Đây là một điểm thiết kế đúng hướng vì cùng một codebase có thể chạy local, Docker hoặc VPS chỉ bằng cách thay đổi cấu hình.

\chapter{QUY TRÌNH VẬN HÀNH VÀ CÁC LUỒNG DỮ LIỆU}
\section{Luồng thiết lập và ủy quyền thiết bị lần đầu}
Quy trình vận hành bắt đầu từ giai đoạn người dùng đưa thiết bị ESP32-S3 vào mạng Wi-Fi gia đình. ESP32 khởi động BLE provisioning với service name \texttt{Prov\_ESP32\_Health}. Ứng dụng Android quét BLE, phát hiện thiết bị có tiền tố \texttt{Prov\_}, tạo secure session bằng Espressif Security 1 và gửi SSID/mật khẩu Wi-Fi qua kênh đã mã hóa. Khi nhận credential thành công, ESP32 lưu cấu hình vào NVS, kết nối Wi-Fi và đồng bộ thời gian NTP để timestamp dữ liệu đo có ý nghĩa.

Sau provisioning mạng, hệ thống cloud thực hiện luồng ủy quyền thiết bị. ESP32 gửi MAC và session ID tới backend để nhận device code và user code. Backend lưu phiên vào Redis với TTL 300 giây và mô phỏng ghi thông tin này lên Firebase RTDB để app có thể hiển thị/nhận mã. Người dùng đã đăng nhập trên app duyệt thiết bị bằng JWT và chữ ký HMAC PoP. ESP32 polling token endpoint; khi phiên được duyệt, backend phát hành device access token, lưu quan hệ MAC--user vào MongoDB và xóa phiên tạm khỏi Redis.

\section{Luồng nhận và ghi dữ liệu sinh hiệu hằng ngày}
Ở bản firmware hiện tại, ESP32 gửi dữ liệu trực tiếp lên Firebase RTDB tại \texttt{/devices/\{mac\}/latest}. Khi phiên đo kết thúc, dữ liệu cũng được ghi vào \texttt{/devices/\{mac\}/history/\{timestamp\}}. Android app lắng nghe đường dẫn latest bằng Firebase listener, nhờ đó màn hình có thể cập nhật gần thời gian thực mà không cần polling liên tục.

Ở kiến trúc cloud mở rộng, telemetry sẽ đi theo tuyến ESP32 $\rightarrow$ EMQX MQTT Broker $\rightarrow$ Go MQTT Worker $\rightarrow$ MongoDB Bucket $\rightarrow$ Firebase RTDB $\rightarrow$ Android App. Payload MQTT dạng \texttt{75,98} được worker parse thành BPM và SpO2, sau đó đánh giá trạng thái Normal/Warning và đẩy vào bucket lịch sử theo giờ. Cách tách này giúp backend trở thành điểm kiểm soát bảo mật trung tâm thay vì để thiết bị ghi trực tiếp vào Firebase.

\section{Luồng tư vấn AI và hồ sơ sức khỏe}
Song song với dữ liệu IoT, ứng dụng Android lưu hồ sơ sức khỏe, lịch sử hội thoại và symptom cache. Khi người dùng hỏi AI, app gửi prompt kèm ngữ cảnh tóm tắt như bệnh sử, triệu chứng đã ghi nhận và cơ sở y tế gần vị trí người dùng nếu có. Cơ chế này giúp phản hồi có tính liên tục hơn và tránh hỏi lặp lại thông tin đã biết. Tuy nhiên, phản hồi AI chỉ mang tính tham khảo, không thay thế chẩn đoán y khoa.
\chapter{THẢO LUẬN VÀ PHÂN TÍCH BẢO MẬT}
\section{Kiểm thử backend theo hướng TDD}
Backend có hai file test chính: \texttt{auth\_test.go} và \texttt{device\_test.go}. Các test không phụ thuộc MongoDB hoặc Redis thật mà dùng mock \texttt{DatabaseService}. Cách làm này phù hợp với TDD vì handler có thể được kiểm thử nhanh, cô lập và bao phủ cả nhánh thành công lẫn nhánh lỗi.

\begin{table}[H]
\centering
\caption{Các test case tiêu biểu của backend}
\renewcommand{\arraystretch}{1.15}
\setlength{\tabcolsep}{6pt}
\begin{tabularx}{\textwidth}{p{7.5cm}X}
\toprule
\textbf{Test case} & \textbf{Nội dung kiểm thử} \\
\midrule
\texttt{TestRegister\_Success} & Đăng ký thành công và xác nhận mật khẩu được băm đúng bằng BCrypt. \\
\texttt{TestRegister\_DuplicatePhone} & Từ chối đăng ký khi số điện thoại đã tồn tại. \\
\texttt{TestLogin\_Success} & Đăng nhập đúng mật khẩu và nhận JWT. \\
\texttt{TestLogin\_WrongPassword} & Từ chối mật khẩu sai, không phát hành token. \\
\texttt{TestAuthorize\_Success} & Sinh device code 32 ký tự và user code dạng \texttt{AAAA-1111}. \\
\texttt{TestConfirm\_InvalidSignature} & Từ chối phiên có chữ ký PoP sai. \\
\texttt{TestConfirm\_Success} & Duyệt phiên đúng JWT, MAC, session và HMAC. \\
\texttt{TestToken\_Pending} & ESP32 polling trước khi user duyệt nhận lỗi \texttt{authorization\_pending}. \\
\texttt{TestToken\_Success} & Sau khi duyệt, backend sinh access token, lưu device và xóa Redis session. \\
\bottomrule
\end{tabularx}
\end{table}

\section{Giao diện giám sát thời gian thực trên ứng dụng}
Màn hình IoT trên Android đóng vai trò như một dashboard vận hành. Người dùng có thể thấy trạng thái kết nối BLE, thiết bị đang được phát hiện, trạng thái provisioning, nhịp tim, SpO2, nhiệt độ, độ ẩm và trạng thái phiên đo. Dữ liệu được cập nhật bằng Firebase listener nên độ trễ hiển thị phụ thuộc chủ yếu vào đường truyền Internet và tốc độ ghi RTDB của ESP32.

\begin{figure}[H]
    \centering
    \framebox[0.9\textwidth]{\raisebox{0pt}[6.0cm][6.0cm]{
        \begin{minipage}{0.82\textwidth}
            \centering
            \textbf{[Bố cục dashboard IoT trên Android]} \\
            \vspace{0.35cm}
            \begin{tabular}{|p{7cm}|p{6cm}|}
            \hline
            \textbf{Thông số sinh hiệu} & \textbf{Trạng thái thiết bị} \\
            Nhịp tim hiện tại: 75 BPM & BLE: Connected / Disconnected \\
            SpO2: 98\% & Provisioning: CONNECTING / SUCCESS \\
            Nhiệt độ: 28.5$^{\circ}$C & Firebase path: devices/\{MAC\}/latest \\
            Độ ẩm: 72\% & Session: IDLE / MEASURING / COMPLETED \\
            Xu hướng 20 điểm BPM gần nhất & hasFinger: true / false \\
            \hline
            \end{tabular}
        \end{minipage}
    }}
    \caption{Bố cục logic của giao diện giám sát dữ liệu sinh hiệu thời gian thực.}
\end{figure}

\section{Bảo mật giao tiếp ban đầu giữa ứng dụng và thiết bị}

Giai đoạn giao tiếp ban đầu giữa ứng dụng Android và ESP32-S3 là một trong những điểm nhạy cảm nhất của hệ thống. Khi thiết bị chưa có cấu hình Wi-Fi, app phải kết nối tới thiết bị qua BLE để truyền SSID và mật khẩu mạng. Nếu quá trình này được thực hiện bằng plaintext hoặc chỉ dựa vào tên thiết bị BLE, attacker ở gần có thể nghe lén gói tin, dựng thiết bị giả có cùng tiền tố tên hoặc chiếm quyền cấu hình thiết bị.

Trong mã nguồn hiện tại, hệ thống sử dụng Espressif BLE Provisioning ở mức Security 1. Ở phía Android, ứng dụng tạo thiết bị ESP qua kênh BLE và cấu hình mức bảo mật tương ứng. Ở phía firmware, ESP32-S3 cũng khởi động Wi-Fi provisioning với cùng mức bảo mật này. Nhờ đó, hệ thống thiết lập secure session trước khi truyền credential, giúp giảm nguy cơ lộ SSID và mật khẩu Wi-Fi trên đường truyền BLE.


Bên cạnh đó, hệ thống sử dụng Proof of Possession (PoP) để chứng minh app đang cấu hình đúng thiết bị. Tuy nhiên, PoP hiện đang là giá trị tĩnh 12345678 trong cả firmware và Android app, phù hợp cho demo nhưng chưa đủ an toàn khi triển khai thật. Trong phiên bản production, mỗi thiết bị nên có PoP riêng, được in dưới dạng QR code hoặc sinh trong quá trình sản xuất. App cũng nên xác thực thêm MAC, service UUID hoặc session ID để tránh kết nối nhầm thiết bị giả mạo.

Sau khi provisioning thành công, ESP32 nên ngắt kết nối BLE và dừng chế độ provisioning nếu không còn cần thiết. Việc thu hẹp thời gian BLE mở giúp giảm bề mặt tấn công cục bộ và ngăn attacker tiếp tục thử kết nối lại vào thiết bị sau khi đã được đưa vào mạng Wi-Fi.

\section{Mô hình đe dọa theo từng tầng}
Đối với hệ thống IoT y tế, rủi ro bảo mật cần được phân tích theo từng tầng thay vì chỉ đặt HTTPS ở backend. Bảng dưới đây tổng hợp các tài sản cần bảo vệ, mối đe dọa và biện pháp giảm thiểu trong thiết kế hiện tại.

\begin{longtable}{
p{0.1\textwidth}
p{0.22\textwidth}
p{0.25\textwidth}
p{0.3\textwidth}
}
\caption{Threat model của hệ thống IoT y tế}\\
\toprule
\textbf{Tầng} & \textbf{Tài sản cần bảo vệ} & \textbf{Mối đe dọa} & \textbf{Biện pháp giảm thiểu} \\
\midrule
\endfirsthead

\toprule
\textbf{Tầng} & \textbf{Tài sản cần bảo vệ} & \textbf{Mối đe dọa} & \textbf{Biện pháp giảm thiểu} \\
\midrule
\endhead

BLE provisioning & SSID, mật khẩu Wi-Fi, quyền cấu hình thiết bị & Nghe lén BLE, cấu hình nhầm thiết bị, MITM local & Espressif Security 1, PoP, khuyến nghị PoP riêng từng thiết bị. \\
ESP32 firmware & Firebase key, định danh thiết bị, dữ liệu cảm biến & Trích xuất firmware, giả mạo MAC, reset Wi-Fi trái phép & Tách secret ra file private, dùng MAC làm device ID, hướng tới token thiết bị. \\
HTTP API & Tài khoản người dùng, token, phiên ghép cặp & Brute-force login, replay user code, chiếm quyền pair & BCrypt, JWT, Redis TTL, HMAC PoP, kiểm tra MAC/session. \\
MQTT broker & Telemetry realtime, topic của từng thiết bị & Publish giả topic, nghe lén telemetry, flood message & Đề xuất MQTT over TLS, ACL theo MAC và device token. \\
Database & Lịch sử sinh hiệu, device token, user profile & Lộ dữ liệu y tế, truy vấn vượt quyền, log nhạy cảm & MongoDB auth, Firebase rules, ẩn token trong log, mã hóa at-rest. \\
Android app & Hồ sơ sức khỏe, lịch sử chat, vị trí & Lộ dữ liệu local, lạm dụng AI context, quyền location & Room local, giới hạn context gửi AI, xin quyền theo nhu cầu. \\

\bottomrule
\end{longtable}

\begin{landscape}
% longtable ở đây
\end{landscape}

\section{Đánh giá những điểm đã làm tốt và điểm còn yếu}
Hệ thống có một số điểm thiết kế tốt so với nguyên mẫu IoT cơ bản. Thứ nhất, provisioning không hard-code Wi-Fi mà dùng BLE secure session. Thứ hai, backend không gắn trực tiếp thiết bị với tài khoản bằng MAC đơn thuần mà có Device Flow với Redis TTL và HMAC PoP. Thứ ba, dữ liệu lịch sử không ghi thô từng điểm đo thành document riêng mà đã có Hourly Bucket Pattern. Thứ tư, code backend có interface database nên unit test có thể mock dễ dàng.

Tuy nhiên, vẫn còn các điểm cần cải thiện trước khi triển khai thật. Firmware hiện tại upload trực tiếp Firebase RTDB bằng legacy token, chưa publish MQTT qua TLS tới EMQX. Backend mới mô phỏng Firebase Admin SDK bằng log. PoP và JWT secret đang hard-code. MQTT trong Docker Compose đang mở port 1883 plaintext, chưa có ACL. App Android còn đường dẫn thiết bị mặc định \texttt{1020BA49D1C8}, cần gắn với kết quả pairing backend thay vì cấu hình tĩnh.

\chapter{KẾT LUẬN VÀ HƯỚNG PHÁT TRIỂN}
\section{Kết luận}
Đề tài đã xây dựng một nguyên mẫu hệ thống IoT y tế có đầy đủ các tầng quan trọng: thiết bị ESP32-S3 đọc cảm biến sinh hiệu, ứng dụng Android cấu hình và hiển thị dữ liệu, backend Golang quản lý xác thực/ủy quyền thiết bị, Redis lưu phiên tạm, MongoDB lưu lịch sử, EMQX làm broker MQTT và Firebase RTDB phục vụ realtime. So với một project đọc cảm biến đơn giản, hệ thống này có giá trị học thuật ở chỗ mô phỏng được nhiều vấn đề thật của IoT: provisioning an toàn, identity của thiết bị, ghép cặp thiết bị với người dùng, truyền dữ liệu realtime, lưu telemetry tối ưu và phân tích threat model.

Về mặt hiện thực, firmware đã có máy trạng thái đo 15 giây, đồng bộ dữ liệu bằng mutex và upload Firebase. Android app có BLE provisioning, StateFlow, Firebase listener, Room và Gemini AI. Backend có API đăng ký/đăng nhập, BCrypt, JWT, Device Flow, HMAC PoP, MQTT worker và unit test. Các thành phần này cho thấy nhóm đã không chỉ viết giao diện minh họa mà đã xây dựng được một kiến trúc có thể mở rộng thành sản phẩm IoT hoàn chỉnh.

Dù vậy, hệ thống vẫn là nguyên mẫu học phần. Các cơ chế như MQTT over TLS, ACL EMQX, Firebase Admin SDK, quản lý secret, rate limiting và token thiết bị riêng cần được hoàn thiện để đạt mức an toàn cao hơn. Ngoài ra, kết quả đo BPM/SpO2 từ MAX30102 trong project chỉ nên xem là tham khảo kỹ thuật, chưa phải kết quả y tế có hiệu lực chẩn đoán.

\section{Hướng phát triển}
Hướng phát triển quan trọng nhất là thống nhất luồng dữ liệu cloud. Firmware cần được chuyển từ upload trực tiếp Firebase sang publish MQTT qua TLS tới EMQX. Backend sẽ xác thực device token, kiểm tra ACL topic, ghi MongoDB bucket và cập nhật Firebase RTDB bằng Admin SDK. Khi đó Firebase chỉ đóng vai trò kênh realtime cho app, còn backend là điểm kiểm soát bảo mật trung tâm.

Hướng thứ hai là nâng cấp bảo mật định danh thiết bị. Mỗi ESP32 nên có PoP riêng, device token riêng và quy trình revoke token khi thiết bị bị mất. JWT secret, Firebase key và MQTT credential cần được đưa vào secret manager hoặc biến môi trường, không hard-code trong mã nguồn. Backend cũng cần bổ sung rate limiting, khóa tạm user code sau nhiều lần nhập sai và audit log cho các hành động ghép cặp.

Hướng thứ ba là nâng cấp phân tích sức khỏe. Khi MongoDB đã có lịch sử BPM/SpO2 theo bucket, hệ thống có thể phát triển mô hình cảnh báo sớm dựa trên xu hướng thay vì chỉ so ngưỡng tức thời. App Android có thể kết hợp dữ liệu sinh hiệu, hồ sơ bệnh nền, triệu chứng người dùng nhập và vị trí địa lý để đưa ra khuyến nghị an toàn hơn, ví dụ gợi ý đi khám khi SpO2 giảm kéo dài hoặc nhịp tim bất thường lặp lại nhiều phiên đo.

\newpage
\appendix
\chapter{PHỤ LỤC: CẤU TRÚC MÃ NGUỒN}
\begin{lstlisting}
IOT/
  front-end/
    app/src/main/java/edu/hust/medicalaichatbot/
      data/                 # Room, Firebase AI, BLE, Location, repositories
      domain/               # Model va use case
      ui/                   # Compose screens, components, viewmodels
      utils/                # Constants, Markdown, PreferenceManager
  medical-ai-chatbot-backend/
    main.go                 # Khoi tao MongoDB, Redis, MQTT worker, Gin router
    auth_handler.go         # Register/Login, BCrypt, JWT
    device_handler.go       # OAuth 2.0 Device Authorization Flow
    mqtt_worker.go          # Subscribe devices/+/telemetry
    database.go             # DatabaseService, MongoDB, Redis
    models.go               # User, Device, TelemetryHistory, payloads
    auth_test.go            # Unit test authentication
    device_test.go          # Unit test Device Flow
    docker-compose.yml      # App, MongoDB, Redis, EMQX
  medical-ai-chatbot-firmware/
    platformio.ini
    src/main.cpp            # ESP32-S3 firmware
\end{lstlisting}

\chapter{PHỤ LỤC: LỆNH CHẠY THAM KHẢO}
\begin{lstlisting}[language=bash]
# Backend
cd medical-ai-chatbot-backend
docker compose up -d --build
go test -v ./...

# Firmware
cd medical-ai-chatbot-firmware
pio run
pio run --target upload
pio device monitor

# Android
cd front-end
./gradlew assembleDebug
\end{lstlisting}

\end{document}




